\chapter{Traffic Congestion}
\label{chap:traffic-congestion}

This chapter presents the development of the traffic congestion module, which addresses one of Ho Chi Minh City's most pressing urban challenges: endless traffic congestion, especially during rush hours. To help commuters avoid traffic jams, we propose a traffic-aware routing feature that leverages real-time traffic camera feeds to analyze current traffic conditions and suggest alternative routes.

The system processes live traffic data from multiple camera sources to identify congestion hotspots, assess severity levels, and provide actionable insights for route optimization and informed travel decisions.
\newline\textit{The underlying camera ingestion and preprocessing pipeline is documented in detail in Chapter~\ref{chap:data-collection} (Data Collection).}



\section{Problem Definition}

Ho Chi Minh City experiences severe and persistent traffic congestion, particularly during peak commuting hours. This congestion:
\begin{itemize}
    \item Causes significant delays and unpredictable travel times
    \item Reduces overall mobility and economic productivity
    \item Creates frustration for daily commuters and visitors
    \item Lacks real-time visibility for route planning decisions
\end{itemize}

\section{Decomposition}
\forestset{
    overviewnode/.style={text width=4cm, font=\small},
    mini/.style={text width=3cm, font=\footnotesize}
}

\begin{figure}[H]
    \centering
    \resizebox{\linewidth}{!}{%
    \begin{forest}
        for tree={
            grow'=south,
            rounded corners,
            draw=black!70, % Default border color
            line width=0.8pt,
            node options={align=left, inner xsep=5pt, inner ysep=4pt},
            text width=4.4cm,
            s sep=7mm,
            l sep=9mm,
            anchor=north,
            edge={-stealth, line width=0.8pt, draw=black!50},
            fill=white % Default fill
        }
            % 6. TRAFFIC CONGESTION (Red theme)
            [
                {Traffic Congestion}, overviewnode, for tree={fill=red!10, draw=red!60},
                [
                    {Vehicle Detection System}, overviewnode,
                    [{YOLOv11 Integration}, mini]
                    [{Vehicle Counting}, mini]
                ]
                [
                    {Data Aggregation Layer}, overviewnode,
                    [{Congestion Level Classification}, mini]
                    [{Spatial Aggregation}, mini]
                ]
                [
                    {REST API (Flask)}, overviewnode,
                    [{Congestion Data Endpoints}, mini]
                    [{Real-time Updates}, mini]
                ]
                [
                    {Web Frontend}, overviewnode,
                    [{Traffic Visualization}, mini]
                    [{Route Suggestion Interface}, mini]
                ]
            ]
    \end{forest}
    }
    \caption{Traffic Congestion Module Decomposition}
    \label{fig:traffic-congestion-overview}
\end{figure}

To address these challenges, we developed a \textbf{traffic-aware routing feature} that:
\begin{itemize}
    \item Leverages real-time traffic camera feeds from across the city
    \item Analyzes current traffic conditions using computer vision
    \item Detects vehicle density and traffic flow patterns
    \item Suggests alternative routes to avoid congested areas
    \item Provides up-to-date traffic information to commuters
\end{itemize}

\section{Framework \& Tech Stack}

The traffic congestion system follows a multi-stage pipeline architecture:

\begin{enumerate}
    \item \textbf{Data Collection Pipeline:} Multi-threaded crawler continuously fetches traffic camera images and stores them in MySQL database
    
    \item \textbf{Vehicle Detection System:} YOLOv11-based computer vision model analyzes images to detect and count vehicles
    
    \item \textbf{Data Aggregation Layer:} Processes detection results to compute congestion metrics and traffic density levels
    
    \item \textbf{REST API (Flask):} Provides access to processed traffic data and congestion information for client applications
    
    \item \textbf{Web Frontend:} Displays real-time traffic conditions and integrates with routing features
\end{enumerate}
    
\section{Implemented Tasks}

\subsection{Vehicle Detection System}

\begin{itemize}
    \item Integrated \textbf{YOLOv11} object detection model
    \item Pretrained model for detecting various vehicle types (cars, motorcycles, buses, trucks)
    \item Achieved high detection accuracy suitable for real-time traffic analysis
    \item Implemented batch processing capabilities for efficient image analysis
    \item Optimized inference speed to handle high-frequency camera feeds
    \item Built vehicle counting and tracking logic to assess traffic density
\end{itemize}

\subsection{Data Aggregation Layer}

\begin{itemize}
    \item Developed algorithms to aggregate vehicle detection results across multiple camera locations
    \item Implemented \textbf{congestion level classification} based on vehicle density thresholds
    \item Created time-series data structures for tracking congestion patterns over time
    \item Built spatial aggregation to compute area-wide traffic conditions
    \item Implemented statistical analysis for identifying traffic trends and anomalies
    \item Designed efficient data structures for fast query and retrieval
\end{itemize}

\subsection{REST API Development}

\textbf{Completed Components:}
\begin{itemize}
    \item Built \textbf{Flask-based REST API} endpoints for accessing traffic data
    \item Implemented endpoints for current congestion status by location
    \item Created API for historical traffic pattern queries
    \item Developed endpoints for camera image retrieval and metadata
    \item Implemented authentication and rate limiting for API security
    \item Created comprehensive API response formatting with JSON schemas
\end{itemize}

\textbf{Ongoing Work:}
\begin{itemize}
    \item \textbf{Testing integration with frontend} to ensure seamless data flow
    \item \textbf{Building real-time server container} for the YOLOv11 model to enable live processing
    \item Optimizing API response times for high-traffic scenarios
    \item Implementing WebSocket support for push-based real-time updates
\end{itemize}

\subsection{Web Frontend}

\textbf{Completed Components:}
\begin{itemize}
    \item Developed base UI components for traffic visualization
    \item Implemented map integration displaying traffic camera locations
    \item Created traffic condition overlay with color-coded severity indicators
    \item Built basic route suggestion interface with congestion awareness
    \item Implemented user preference settings for traffic alerts
\end{itemize}

\textbf{Ongoing Work:}
\begin{itemize}
    \item Continuous editing and refinement of user interface elements
    \item Enhancing real-time update mechanisms for live traffic data
    \item Improving responsive design for mobile device compatibility
    \item Adding interactive features for exploring traffic patterns
    \item Integrating with routing module for comprehensive route planning
\end{itemize}

\subsection{Technology Stack}

\begin{itemize}
    \item \textbf{Programming Language:} Python
    \item \textbf{Web Framework:} Flask (REST API server)
    \item \textbf{Computer Vision:} YOLOv11 (vehicle detection)
    \item \textbf{Database:} MySQL (image and metadata storage)
    \item \textbf{Concurrency:} Multi-threading (parallel data collection)
    \item \textbf{Frontend:} React (web interface)
    \item \textbf{Containerization:} Docker (model deployment)
\end{itemize}

\section{Next Steps}

Future development priorities include:

\begin{itemize}
    \item \textbf{Complete REST API Integration:}
    \begin{itemize}
        \item Finalize frontend integration testing
        \item Deploy production-ready real-time server container
        \item Implement comprehensive error handling and monitoring
    \end{itemize}
    
    \item \textbf{Complete Web Frontend:}
    \begin{itemize}
        \item Polish user interface and user experience
        \item Complete mobile responsiveness optimization
        \item Finalize integration with routing and map features
    \end{itemize}
    
    \item \textbf{Predictive Analytics:} Implement machine learning models to forecast traffic conditions 15-60 minutes ahead based on historical patterns
    
    \item \textbf{Incident Detection:} Develop anomaly detection to identify accidents, breakdowns, or unusual congestion patterns automatically
    
    \item \textbf{Performance Optimization:} Optimize model inference speed and API response times for scalability to city-wide deployment
    
    \item \textbf{Alert System:} Build notification service to alert users about sudden traffic changes or congestion on saved routes
    
    \item \textbf{Historical Analysis Dashboard:} Create analytics tools for studying long-term traffic trends and patterns
    
    \item \textbf{Multi-camera Fusion:} Enhance spatial aggregation to provide comprehensive traffic view across connected road segments
\end{itemize}

\section{Summary}

The traffic congestion module has made substantial progress in addressing Ho Chi Minh City's traffic challenges through a computer vision-driven architecture. It builds on the multi-threaded ingestion foundation described in Chapter~\ref{chap:data-collection}, combining: (1) the completed camera data pipeline (100\%), (2) the YOLOv11 vehicle detection system (100\%), and (3) the aggregation layer (100\%). The REST API is 70\% complete (focus: real-time containerization and frontend integration) and the web frontend stands at 60\% (UI refinement and interactive exploration). This traffic-aware system converts raw camera frames into actionable congestion intelligence for routing decisions and user situational awareness, positioning the platform to improve urban mobility through proactive avoidance of emerging hotspots.
